
\def\PUDcodigo{ACE017}

\if\COMPILAR0
    \def\PUDcodacad{04507.5}
\else
    \def\PUDcodacad{04506.5}
\fi

\def\PUDdisciplina{Introdução à Engenharia}

\def\PUDsemestre{S1}

\def\PUDhoras{40}

%NOVOS ---------------------------------------------------
\def\PUDhorasTeoria{40}
\def\PUDhorasPratica{0}

\if\COMPILAR0
   \def\PUDinterECA{
   \hyperref[PUD:ACE007]{Ética e Cidadania (04507.33)}
   
   \hyperref[PUD:ACE015]{Émpreendedorismo (04507.38)}
   }
\else
   \def\PUDinterECA{
   \hyperref[PUD:ACE007]{Ética e Cidadania (04506.32)}
   
   \hyperref[PUD:ACE015]{Émpreendedorismo (04506.37)}
   }
\fi

\def\PUDatuaisECA{
- Solucionando problemas com criatividade por meio de metodologias como ``Pensamento Lateral'', ``Teoria da Resolução Inventiva de Problemas (TRIZ)'', \textit{Design Thinkin}, etc.

- Introdução à Indústria 4.0
}

\def\PUDrecursos{Projetor multimídia; Quadro branco e pincel.}

%----------------------------------------------------------

\def\PUDcreditos{2}

\def\PUDprerequisito{---}

\def\PUDpreacad{---}

\def\PUDobjetivos{Conhecer a instituição de ensino e os objetivos do curso de engenharia.  Compreender pontos sobre a atuação e responsabilidades dos profissionais de engenharia.  Conhecer ferramentas básicas de planejamento e simulação para engenharia. Conhecer os conceitos básicos da ciência econômica e da engenharia econômica discutindo os aspectos relacionados ao setor industrial. Conhecer metodologias para solução de problemas.}

\def\PUDementa{Regulamento e normas para o ensino no IFCE, Apoio institucional do IFCE ao discente, Direitos e deveres do aluno. Histórico da Engenharia, Projeto em engenharia, modelos e simulação, legislação profissional do engenheiro, sistema CONFEA/CREAs. Organização didática do curso de Engenharia, estruturação do curso em suas áreas, campos de atuação do engenheiro, pesquisa tecnológica e pesquisa científica, comunicação em engenharia nas formas escrita, gráfica e oral, perfil do engenheiro, conhecimento de idiomas, habilidade empreendedora, responsabilidade social e conduta ética. Teoria geral sobre economia. Engenharia econômica. Tópicos de economia aplicados a engenharia. Solução de problemas.}

\def\PUDconteudo{ & O ensino no IFCE: Direitos e Deveres do aluno. Oficina de acolhimento proporcionada pelos diversos setores da administração do campus.
\\ 
 &  História da engenharia. Legislação profissional, CONFEA e CREAs. Organização do \NOMECURSO do IFCE. 
\\ 
& Visação Geral do Mercado de Trabalho do \NOMEEGRESSO. Introdução à Indústria 4.0.
\\
 & O Engenheiro: O que é um Engenheiro? Engenheiro – Aquele que Busca Soluções. Demanda por Engenheiros. Equipe Tecnológica. Método de Projeto de Engenharia. Especialidades da Engenharia e Áreas Afins. Funções da Engenharia. Engenheiro como um Profissional. Características de um Engenheiro de Sucesso. Características de um Engenheiro Criativo.
 %Mercado de trabalho.  Perfil do engenheiro. 
\\ 
& Solucionando Problemas: Tipos de Problemas. Procedimentos para Solucionar Problemas. Aptidões para Solucionar Problemas. Técnicas para a Solução de Problemas sem Erros. Estimativas. Soluções Criativas de Problemas.
\\
 & Desenvolvimento de um projeto de engenharia e Análise de investimentos.  O que é Pesquisa. Formatação de livros usando o Word.  Equações em planilhas eletrônicas.  Gráficos em planilhas eletrônicas.  Digitação de fórmulas em editores de texto.  Formatação de planilhas.  %Ferramenta computacional Matlab.
\\
 & Introdução à Economia.
 %A ciência econômica. Microeconomia: oferta, demanda e mercado. Elasticidade e estrutura de mercado (concorrência perfeita, oligopólio e monopólio). Macroeconomia: teoria geral do emprego,juros e moeda. O sistema financeiro e o Banco Central. Inflação, recessão e endividamento. Economia Brasileira: desenvolvimento e subdesenvolvimento. 
 }

\def\PUDmetodo{Aulas expositórias ministradas pelo professor, apresentações dos diversos setores do IFCE feitas por seus próprios servidores, palestras realizadas por engenheiros da área e profissionais especialistas em certos temas específicos.}

\def\PUDavalia{A avaliação será feita com base na participação do aluno nas atividades propostas ao longo da disciplina, como apresentações, pequenos projetos de engenharia, trabalhos, relatório, dentre outras.}

\def\PUDbibbasica{ & \href{www.biblioteca.ifce.edu.br/index.asp?codigo_sophia=39514}{Holtzapple}, Mark Thomas.  Introdução à engenharia.  Rio de Janeiro, RJ: LTC, 2013.  220p.  ISBN: 9788521615118. 
\\ 
 &  \href{www.biblioteca.ifce.edu.br/index.asp?codigo_sophia=62399}{Sá}, Antônio Lopes de.  Ética profissional.  9º ed.  São Paulo, SP : Atlas, 2015.  312p.  ISBN: 9788522455348.
% Instituto Federal de Educação, Ciência e Tecnologia do Ceará.  Projeto Pedagógico Curso de Engenharia Controle e Automação.  Maracanaú: IFCE, 2013.  p.  181.
\\ 
 &  \href{www.biblioteca.ifce.edu.br/index.asp?codigo_sophia=61932}{Srour}, Robert Henry.  Ética empresarial.  4º ed.  Rio de Janeiro, RJ : Elsevier, 2013.  213p.  ISBN: 9788535264470.
% Instituto Federal de Educação, Ciência e Tecnologia do Ceará.  Regulamento da Organização Didática.  Fortaleza: IFCE, 2010.  p.  64.  Disponível em: \url{http://www. ifce. edu. br/images/stories/menu_superior/Ensino/ROD/ROD-Comisso_de_Sistematizao27. pdf}.  Data: 18/12/2013. 
}

%{ & Instituto Federal de Educação, Ciência e Tecnologia do Ceará.  Projeto Pedagógico Curso de Engenharia Mecânica.  Maracanaú: IFCE, 2013.  p.  181. 
%\\ 
% & Instituto Federal de Educação, Ciência e Tecnologia do Ceará.  Projeto Pedagógico Curso de Engenharia Controle e Automação.  Maracanaú: IFCE, 2013.  p.  181.   
%\\ 
% & Instituto Federal de Educação, Ciência e Tecnologia do Ceará.  Regulamento da Organização Didática.  Fortaleza: IFCE, 2010.  p.  64.  Disponível em: \url{http://www. ifce. edu. br/images/stories/%menu_superior/Ensino/ROD/ROD-Comisso_de_Sistematizao27. pdf}.  Data: 18/12/2013. }

\def\PUDbibcomplementar{ 
& \href{www.biblioteca.ifce.edu.br/index.asp?codigo_sophia=40759}{Motta}, Ronaldo Seroa da.  Economia ambiental.  Rio de Janeiro, RJ : FGV, 2013.  225p.  ISBN: 8522505446.
\\
& \href{biblioteca.ifce.edu.br/index.asp?codigo_sophia=65379}{Salim}, Cesar Simões; SILVA, Nelson Caldas. Introdução ao empreendedorismo: despertando a atitude empreendedora. Rio de Janeiro: Elsevier, 2010. 245 p. (Coleção Empreendedorismo). ISBN 9788535234664.
\\ 
 %&  \href{www.biblioteca.ifce.edu.br/index.asp?codigo_sophia=23755}{Vesilind}, P. Aarne.  Introdução à engenharia ambiental.  São Paulo, SP : Cengage Learning, 2011.  438p.  ISBN: 9788522107186.
% Conselho Regional de Engenharia e Agronomia do Ceará.  Legislação.  Disponível em: http://www. creace. org. br/.  Data: 18/12/2013. 
%\\ 
 %&  \href{www.biblioteca.ifce.edu.br/index.asp?codigo_sophia=24082}{Braga}, Benedito.  Introdução à engenharia ambiental.  2º ed.  São Paulo, SP : Pearson Prentice Hall, 2005.  318p.  ISBN: 9788576050414.
%\\
 &  \href{www.biblioteca.ifce.edu.br/index.asp?codigo_sophia=58171}{Freitas}, Carlos Alberto de.  Introdução à Engenharia.  E-book.  Pearson.  160p.  ISBN: 9788543005515.
\\
 &  \href{www.biblioteca.ifce.edu.br/index.asp?codigo_sophia=24233}{Camargo}, Marculino.  Fundamentos de ética geral e profissional.  10º ed.  Petrópolis, RJ: Vozes, 2011.  108p.  ISBN: 9788532621313. 
 \\
 & \href{http://ifce.bv3.digitalpages.com.br/users/publications/9788576053590}{Samanez}, Carlos Patrício. Engenharia Econômica. E-book: Pearson. 226 p. ISBN 9788576053590.
 %\\
 %& \href{http://ifce.bv3.digitalpages.com.br/users/publications/9788582124000 }{Lenzi}, Ervin K.; Lenzi, MARCELO K.; Ryba, Andrea. Elementos de Engenharia Econômica. E-book: InterSaberes. 0 p. ISBN 9788582124000.
 }

\def\PUDautor{Geraldo Luis Bezerra Ramalho}

\def\PUDdata{2013-04-17}

\def\PUDrevisao{3}

\def\PUDrevisor{Celso Rogério Schmidlin Júnior}

\def\PUDdatarevisao{2019-05-15}

\PUDcorpo
